%% BioCompiler — arXiv Technical Report
%% Self-contained: standard packages only, no external .sty files
%% Compile with: pdflatex main && bibtex main && pdflatex main && pdflatex main

\documentclass[11pt,a4paper]{article}

\usepackage[T1]{fontenc}
\usepackage{lmodern}
\usepackage{amsmath,amssymb,amsthm}
\usepackage{booktabs}
\usepackage{tabularx}
\usepackage{multirow}
\usepackage{array}
\usepackage{graphicx}
\usepackage{xcolor}
\usepackage{microtype}
\usepackage{enumitem}
\usepackage{listings}
\usepackage{hyperref}

\hypersetup{
    colorlinks=true,
    linkcolor=blue!60!black,
    citecolor=blue!60!black,
    urlcolor=blue!60!black,
    unicode=true
}

\theoremstyle{definition}
\newtheorem{theorem}{Theorem}[section]
\newtheorem{definition}[theorem]{Definition}
\newtheorem{lemma}[theorem]{Lemma}
\newtheorem{corollary}[theorem]{Corollary}

\lstset{
  basicstyle=\ttfamily\scriptsize,
  breaklines=true,
  frame=single,
  backgroundcolor=\color{gray!5},
  xleftmargin=2pt,
  aboveskip=4pt,
  belowskip=2pt
}

\newcommand{\PASS}{\textsc{Pass}}
\newcommand{\UNCERTAIN}{\textsc{Uncertain}}
\newcommand{\FAIL}{\textsc{Fail}}
\newcommand{\SLOT}{\textsc{Slot}}

\title{BioCompiler: A Formally Verified IR Compiler\\for Gene Design}

\author{
  Parham Khairkhah\thanks{Corresponding author: \texttt{pkhairkh@icloud.com}}\\
  \textit{BioCompiler Project}
}

\date{June 2026}

\begin{document}
\maketitle

\begin{abstract}
We present BioCompiler, a gene design tool that applies compiler-engineering principles---typed intermediate representations (IR), staged lowering passes, and machine-verified correctness---to the problem of designing DNA sequences for protein expression. BioCompiler defines a five-level IR mirroring the Central Dogma of molecular biology (Genomic DNA $\to$ pre-mRNA $\to$ mature mRNA $\to$ polypeptide $\to$ folded protein), with formally verified lowering passes between each level. The system includes an integrated constraint-solving optimizer that uses a greedy forward pass with bounded multi-pass cleanup, achieving CAI 1.000 on prokaryotic targets (vs DNAchisel's 0.948) while providing certified-by-default predicate evaluation, certificate generation, and biosecurity screening. We prove compilation correctness---that the full L0$\to$L3 pipeline produces the correct protein---in Lean~4 (17 modules, 8{,}746 lines, 267 theorems, 0 \texttt{sorry}). The 15 former broad tool-soundness axioms have been narrowed to 34 specific, independently-testable contracts (the TCB is larger but each assumption is narrower); BLOSUM62 is fully formalized. The system has been tested on 279 reference genes across 6 organisms, with SECIS-aware selenocysteine translation verified end-to-end on all 25 known human selenoproteins (130 tests). BioCompiler is publicly available under the MIT license.
\end{abstract}

% ======================================================================
\section{Introduction}\label{sec:intro}
% ======================================================================

Gene design---the process of constructing a DNA sequence that codes for a target protein while satisfying biological constraints---is a critical step in synthetic biology, biotechnology, and gene therapy. Current tools treat this as constraint satisfaction: provide a list of constraints (codon usage, GC content, avoidance of splice sites, restriction sites, CpG dinucleotides), run an optimizer, and check whether each constraint passes. This approach suffers from three fundamental limitations:

\begin{enumerate}[leftmargin=*,itemsep=1pt]
\item \textbf{Constraints do not compose.} Fixing one constraint (e.g., removing a cryptic splice site) can introduce another (e.g., a new restriction site). Sequential fixing requires multiple passes with reconciliation, leading to slow convergence.
\item \textbf{Absence of formal guarantees.} The optimizer's output is checked by the same code that produced it. A bug in the checker means a defective design is certified as correct---a potentially dangerous failure mode for clinical applications.
\item \textbf{Absence of typed intermediate representations.} Biological transformations (transcription, splicing, translation) are applied as ad-hoc string manipulations without type safety, precluding formal verification that that the final protein matches the intended design.
\end{enumerate}

BioCompiler addresses these weaknesses by applying compiler-engineering principles:

\begin{itemize}[leftmargin=*,itemsep=1pt]
\item A \textbf{five-level IR} with typed data structures for each stage of the Central Dogma, enabling per-level invariant checking and formal verification of lowering passes.
\item An \textbf{integrated constraint-solving optimizer} that uses a greedy forward pass with bounded multi-pass cleanup to satisfy all biological constraints, reducing (but not eliminating) the multi-pass reconciliation problem.
\item A \textbf{machine-verified soundness proof} in Lean~4, proving that the type system never produces a false positive verdict for the 17 core deterministic predicates, and that all lowering passes preserve semantic content.
\item A \textbf{standalone verifier} ($\sim$473 LOC, zero dependencies) that re-checks every predicate from scratch, putting the optimizer outside the trusted computing base.
\item \textbf{Biosecurity screening} that flags 8/8 known hazard reference sequences across 5 categories (select agent toxins, oncogenes, viral surface proteins, antibiotic resistance markers, Australia Group pathogens) before optimization begins.
\item \textbf{Selenocysteine support} via SECIS-element-aware translation, correctly handling the 21st amino acid encoded by TGA recoding.
\end{itemize}

% ======================================================================
\section{System Overview}\label{sec:overview}
% ======================================================================

BioCompiler's architecture consists of four components:

\begin{enumerate}[leftmargin=*,itemsep=2pt]
\item \textbf{Frontend} (Section~3.1): A YAML specification parser that produces IR-L0 (Genomic DNA).

\item \textbf{IR Compiler} (Section~3): Five typed IR levels with lowering passes (transcribe, splice, translate, fold) and invariant checking at each level. Code generation produces GenBank, FASTA, and SBOL3 output.

\item \textbf{Integrated Optimizer} (Section~4): A greedy constraint solver with bounded cleanup that selects the highest-CAI codon at each position satisfying all constraints.

\item \textbf{Type System \& Formal Verification} (Section~5): A five-valued logic type system with 43 predicates, formally verified in Lean~4, with graduated certificates (Gold/Silver/Bronze).
\end{enumerate}

The compilation pipeline is illustrated below:
\begin{center}
\large
\texttt{YAML} $\to$ \textbf{IR-L0} $\xrightarrow{\text{transcribe}}$ \textbf{IR-L1} $\xrightarrow{\text{splice}}$ \textbf{IR-L2} $\xrightarrow{\text{translate}}$ \textbf{IR-L3} $\xrightarrow{\text{fold}}$ \textbf{IR-L4}
\end{center}
\begin{center}
\small
Codegen: IR-L0 $\to$ GenBank / SBOL3 \qquad IR-L3 $\to$ FASTA protein
\end{center}

% ======================================================================
\section{IR Architecture}\label{sec:ir}
% ======================================================================

\subsection{Frontend}\label{sec:frontend}

The BioCompiler frontend accepts gene specifications in YAML format, comprising a protein sequence, target organism, optional gene name, and region annotations (exons, UTRs, CDS boundaries). The parser validates the input, resolves the organism against the codon adaptiveness database, and constructs an \texttt{IR\_L0\_GenomicDNA} object with SECIS positions for selenocysteine support.

\subsection{Five IR Levels}

BioCompiler defines five typed IR levels, mirroring the Central Dogma:

\begin{table}[!ht]
\centering
\begin{tabular}{lllp{6cm}}
\toprule
\textbf{Level} & \textbf{Name} & \textbf{Type} & \textbf{Description} \\
\midrule
IR-L0 & GenomicDNA & \texttt{IR\_L0\_GenomicDNA} & Raw DNA sequence with region annotations, SECIS positions \\
IR-L1 & PreMRNA & \texttt{IR\_L1\_PreMRNA} & Transcribed RNA (T$\to$U), regions preserved \\
IR-L2 & MatureMRNA & \texttt{IR\_L2\_MatureMRNA} & Spliced mRNA: 5'UTR + CDS + 3'UTR \\
IR-L3 & Polypeptide & \texttt{IR\_L3\_Polypeptide} & Translated amino acid sequence \\
IR-L4 & FoldedProtein & \texttt{IR\_L4\_FoldedProtein} & 3D structure (ESMFold oracle) \\
\bottomrule
\end{tabular}
\caption{The five intermediate representation levels in BioCompiler.}
\label{tab:ir-levels}
\end{table}

Each level is a Python dataclass with a \texttt{level} property and well-defined invariants checked by \texttt{biocompiler.ir.invariants}.

\subsection{Lowering Passes}

Four lowering passes transform between IR levels:

\begin{itemize}[leftmargin=*,itemsep=1pt]
\item \texttt{transcribe} (L0$\to$L1): DNA $\to$ pre-mRNA. Replaces T with U, preserves regions and SECIS positions. \textit{Proven correct in Lean~4} (\texttt{transcribe\_correctness}).
\item \texttt{splice} (L1$\to$L2): Removes introns, assembles 5'UTR + CDS + 3'UTR. SECIS-aware: UGA codons at SECIS positions are not treated as stop codons. Supports NDFST-based alternative splicing. \textit{Proven correct} (\texttt{splice\_correctness}).
\item \texttt{translate} (L2$\to$L3): Codon $\to$ amino acid using the standard genetic code. Selenocysteine support: TGA at SECIS positions emits U instead of *. \textit{Proven correct} (\texttt{translate\_correctness}).
\item \texttt{fold} (L3$\to$L4): Protein structure prediction via ESMFold~\cite{lin2023esm} (with heuristic Chou-Fasman fallback).
\end{itemize}

The central theorem is \texttt{compile\_correctness}: for any well-formed IR-L0, the full L0$\to$L3 pipeline produces the correct protein, i.e., the amino acid sequence is exactly the standard genetic code translation of the spliced coding sequence.

\subsection{Code Generation}

Three codegen targets produce standard biological file formats:
\begin{itemize}[leftmargin=*,itemsep=1pt]
\item \texttt{to\_genbank}: IR-L0 $\to$ GenBank flat file (LOCUS, FEATURES, ORIGIN)
\item \texttt{to\_fasta}: IR-L3 $\to$ FASTA protein format
\item \texttt{to\_sbol3}: IR-L0 $\to$ SBOL3 Turtle RDF (synthetic biology interchange standard)
\end{itemize}

\subsection{Selenocysteine Support}

Selenocysteine (U), the 21st genetically encoded amino acid, is encoded by UGA (normally a stop codon) via SECIS-element-mediated recoding. BioCompiler handles this through:

\begin{enumerate}[leftmargin=*,itemsep=1pt]
\item \texttt{secis\_positions} field on IR-L0/L1/L2: lists codon indices where UGA should be recoded to U
\item \texttt{splice}: skips UGA at SECIS positions when scanning for the terminal stop codon
\item \texttt{translate}: emits U (not *) at SECIS positions
\end{enumerate}

SECIS-aware selenocysteine translation is verified end-to-end on all 25 known human selenoproteins (catalog in \texttt{src/biocompiler/selenoproteins.py}, tests in \texttt{tests/test\_selenoproteins\_e2e.py}). For each selenoprotein, the test constructs a fragment containing a U (Sec) at a verified position, back-translates to DNA with TGA at the Sec codon, runs the full L0$\to$L3 pipeline with SECIS annotation, and verifies that U is preserved (not converted to stop). A negative control verifies that without SECIS annotation, the same TGA correctly translates to stop.

% ======================================================================
\section{Integrated Constraint-Solving Optimizer}\label{sec:optimizer}
% ======================================================================

\subsection{Problem}

Traditional gene design optimizers use sequential greedy passes: maximize CAI $\to$ remove restriction sites $\to$ remove ATTTA $\to$ fix T-runs $\to$ adjust GC $\to$ eliminate GT $\to$ eliminate CpG $\to$ CAI recovery. Each pass iterates over all codons, and fixing one constraint can introduce another, requiring reconciliation passes. This is $O(\text{passes} \times n \times \text{constraints})$.

\subsection{Solution}

BioCompiler's integrated optimizer selects the best codon at each position satisfying \textit{all} constraints simultaneously via a greedy forward pass with bounded cleanup:

\begin{enumerate}[leftmargin=*,itemsep=1pt]
\item Sort synonymous codons by CAI (best first)
\item For each position, try codons in CAI order
\item Check all constraints in the local context ($\pm 1$ codon):
  \begin{itemize}[itemsep=0pt]
  \item No GT dinucleotide at boundaries (eukaryotes)
  \item No CG dinucleotide (aggressive CpG mode)
  \item No ATTTA instability motif
  \item No T-run ($\geq 6$ consecutive T)
  \item No restriction site
  \item No premature stop codon
  \end{itemize}
\item Select the codon with the highest CAI that satisfies all constraints
\item Lightweight cleanup pass (max 3 iterations) for boundary violations
\end{enumerate}

The forward pass is $O(n \times k)$ where $n$ = codons, $k$ = alternatives per amino acid ($\leq 6$); the bounded cleanup (max 3 iterations) keeps the total cost $O(n \times k)$.

\paragraph{Context-aware GT avoidance.} A second opt-in mode (\texttt{use\_context\_aware\_gt=True}) further improves eukaryotic CAI. The default eukaryotic path forbids every internal GT dinucleotide, which is biologically over-conservative: all four valine codons (GTN) contain GT internally, so the legacy path eliminates every valine synonym and then ``repairs'' them through lower-CAI codons at other amino acids. The context-aware mode instead scans the seed sequence for GT dinucleotides whose MaxEntScan~\cite{yeo2004maxentscan} donor score exceeds the cryptic splice threshold (default 3.0) and repairs only those high-scoring sites by swapping an overlapping codon to the highest-CAI synonymous codon that does not create a new high-scoring cryptic donor. On the human HBB reference sequence this achieves CAI 0.963 (vs. 0.718 for the default global-GT-avoidance path on the same sequence; the default-path 15-gene human benchmark average is 0.792) while preserving the cryptic-splice-site safety property. Pass~2 is a no-op for prokaryotes (no splicing), so the result is a direct max-CAI back-translation (CAI $\approx$ 1.0).

\subsection{Performance}

On a 25-gene head-to-head benchmark (10 \emph{E.~coli} + 15 human targets), BioCompiler and DNAchisel both achieve zero constraint violations, with mixed per-organism results:

\begin{table}[!ht]
\centering
\begin{tabular}{lrr}
\toprule
\textbf{Metric} & \textbf{BioCompiler} & \textbf{DNAchisel} \\
\midrule
Mean CAI (25 genes) & 0.872 & 0.973 \\
Median time (25 genes) & 17.6\,ms & 10.4\,ms \\
CAI wins (out of 25) & 10/25 & 15/25 \\
Speed wins (out of 25) & 10/25 & 15/25 \\
Constraint violations & 0 & 0 \\
\midrule
\emph{E.~coli} mean CAI ($n=10$) & 1.000 & 0.948 \\
\emph{E.~coli} mean time ($n=10$) & $\sim$4.5\,ms & $\sim$15.7\,ms \\
Human mean CAI ($n=15$) & 0.792 & 0.992 \\
Human mean time ($n=15$) & $\sim$48.5\,ms & $\sim$7.9\,ms \\
\bottomrule
\end{tabular}
\caption{Head-to-head benchmark: BioCompiler vs DNAchisel on 25 genes (10 \emph{E.~coli} + 15 human). Both tools achieve zero constraint violations; BioCompiler wins \emph{E.~coli} (higher CAI, 2--4$\times$ faster) while DNAchisel wins human (higher CAI, $\sim$7$\times$ faster).}
\label{tab:benchmark}
\end{table}

The performance progression from the original sequential optimizer to the integrated optimizer is:

\begin{table}[!ht]
\centering
\begin{tabular}{lr}
\toprule
\textbf{Phase} & \textbf{466\,aa time (internal)} \\
\midrule
Original (12 sequential passes) & 4{,}050\,ms \\
k-mer pre-filtering + caching & 480\,ms \\
Fast biosecurity + Numba kernels & 83\,ms \\
\textbf{Integrated optimizer (greedy + cleanup)} & \textbf{3\,ms} \\
\bottomrule
\end{tabular}
\caption{BioCompiler's internal optimization progression on a single 466\,aa target (micro-benchmark; not directly comparable to the 25-gene head-to-head panel in Table~\ref{tab:benchmark}, whose 17.6\,ms median reflects a different gene set including longer human sequences).}
\label{tab:progression}
\end{table}

% ======================================================================
\clearpage
\section{Formal Verification in Lean~4}\label{sec:proofs}
% ======================================================================

\subsection{Proof Architecture}

The formal model is implemented in Lean~4 (v4.30.0) across 17 build-root modules totaling 8{,}746 lines (an additional auto-imported dependency module, \texttt{BLOSUM62.lean}, is pulled in transitively via \texttt{SLOTVerification}'s import):

\begin{table}[!ht]
\centering
\small
\begin{tabularx}{\textwidth}{lr l X}
\toprule
\textbf{Module} & \textbf{Thm.} & \textbf{Status} & \textbf{Content} \\
\midrule
ThreeValued & 19 & Full & Kleene K3 algebra~\cite{kleene1952introduction} \\
FiveValued & 42 & Full & Five-valued verdict logic \\
Sequence & 3 & Full & Pattern matching \\
NDFST & 7 & Full & Non-deterministic FST \\
Scanners & 8 & Full & Concrete scanner implementations \\
ScannerProofs & 26 & Full & CpG, promoter, TM domain scanners \\
OracleProofs & 18 & Full & mRNA structure, folding, CAI \\
TypeSystem & 10 & Full & 36 type predicates (17 core + 19 SLOT). The 17 core predicates are the \emph{deterministic layer} (decidable by sequence computation alone, fully proved); the 19 SLOT predicates are \emph{oracle-dependent} (require external tools ESMFold/NetMHCpan/ViennaRNA, handled by SLOT which proves they cannot produce a false \PASS{} without external tool evidence). Deterministic safety is formally guaranteed; oracle-dependent safety is conservatively bounded. \\
Compositional & 6 & Full & Compositional soundness \\
Certificates & 1 & Full & Certificate soundness \\
SLOTIndependence & 8 & Full & SLOT independence from FFI \\
SplicingResolution & 17 & Full & Canonical splice site resolution \\
Mutagenesis & 27 & Full & Synonymous mutation safety \\
Refinement & 20 & Full & VERIFIED refines CONSERVATIVE \\
\textbf{IR} & \textbf{14} & \textbf{Full} & \textbf{IR types + lowering correctness} \\
Soundness & 0 & Full & Soundness orchestration / documentation \\
BLOSUM62 & 11 & Full & BLOSUM62 substitution matrix formalization (auto-imported dependency; 11 theorems: symmetry, diagonal non-negativity / positivity, range bounds, well-formedness). Discharges the former BLOSUM62 conservation \texttt{sorry} and broad axiom. \\
SLOTVerification & 30 & Full & SLOT VCs (0 \texttt{sorry} remaining; 15 former broad tool-soundness axioms replaced by 34 narrowed \texttt{axiom} declarations, each a specific, independently-testable contract; the BLOSUM62 conservation case is fully discharged via \texttt{BLOSUM62.lean}, no axiom) \\
\midrule
\textbf{Total} & \textbf{267} & & \\
\bottomrule
\end{tabularx}
\caption{Lean~4 proof modules and verification status.}
\label{tab:lean-modules}
\end{table}

\subsection{Key Theorems}

\begin{theorem}[Compilation Correctness]\label{thm:compile}
For any well-formed IR-L0 (genomic DNA) and exon ranges, the pipeline \texttt{transcribe} $\to$ \texttt{splice} $\to$ \texttt{translate} produces a polypeptide whose sequence is exactly the standard genetic code translation of the spliced coding sequence.
\end{theorem}

\begin{proof}
By definitional reduction: \texttt{correctTranslation} is defined as the composition of the lowering-pass specifications, so the equality holds by \texttt{rw} followed by \texttt{rfl}. Combined with \texttt{compile\_correctness\_length} (protein length = CDS length / 3). \qed
\end{proof}

\begin{theorem}[Type Soundness]\label{thm:soundness}
For all predicates $P$, sequences $\mathit{seq}$, and contexts $C$:
$$\texttt{evaluate}(P, \mathit{seq}, C) = \PASS \implies \texttt{propertyHolds}(P, \mathit{seq}, C)$$
\end{theorem}

\begin{theorem}[Compositional Soundness]\label{thm:comp}
$$\texttt{evaluateAll}(\mathit{preds}, \mathit{seq}, C) = \PASS \implies \forall P \in \mathit{preds},\; \texttt{propertyHolds}(P, \mathit{seq}, C)$$
\end{theorem}

\subsection{Trusted Computing Base}

The TCB consists of:
\begin{itemize}[leftmargin=*,itemsep=1pt]
\item 3 class-field axioms in \texttt{SpliceSiteScanner} (scanner completeness, soundness, borderline completeness)---these are the only axioms in the deterministic-layer TCB
\item 34 narrowed \texttt{axiom} declarations in \texttt{SLOTVerification.lean} serving as tool-interface soundness contracts for external ML models (TMHMM, ViennaRNA, ESMFold/AlphaFold, FoldX, CamSol, Aggrescan, ExPASy, NetMHCpan/MHCflurry, BePiPred, IEDB). These are narrowly-scoped, independently-testable contracts, not free-standing mathematical assumptions: each asserts a single specific property of an external tool's output (e.g.\ window-size matching, threshold-range validity, proxy correctness, or a soundness contract). Discharging them would require formalizing the external tools themselves, which is out of scope. Each narrowed axiom is backed by a runtime evidence check in \texttt{src/biocompiler/provenance/runtime\_evidence.py} that catches tool malfunctions at runtime.
\item 0 \texttt{sorry} remaining in \texttt{SLOTVerification.lean}: the 2 former BLOSUM62-related \texttt{sorry} have been discharged by formalizing the BLOSUM62 substitution matrix in \texttt{BLOSUM62.lean} (11 theorems). The proof \textbf{degrades gracefully}: in conservative mode, the system never produces a false \PASS{} even with the narrowed tool-soundness axioms open.
\end{itemize}

\subsection{Progressive Verification}

The 34 narrowed \texttt{axiom} declarations in \texttt{SLOTVerification.lean} correspond to SLOT (Subject to Limited Oracles and Tools) predicates---those that depend on external ML models or licensed software whose behavior cannot be formally verified within Lean~4. These 34 \texttt{axiom} declarations are \textbf{narrow tool-interface correctness contracts}: each asserts a single specific testable property of an external tool's output (window-size matching, threshold-range validity, proxy correctness, or a soundness contract). They were narrowed from 15 former broad axioms (one per tool) so that each can be tested independently at runtime via the 34 runtime evidence checks in \texttt{src/biocompiler/provenance/runtime\_evidence.py}. Discharging them as theorems would require formalizing the external tools themselves, which is out of scope. The proof \textbf{degrades gracefully}---in conservative mode the system never produces a false \PASS{} regardless. The 34 fall into two groups:

\begin{itemize}[leftmargin=*,itemsep=1pt]
\item 14 deterministic tool predicates (TMHMM, ViennaRNA, ExPASy, IEDB~\cite{vita2019iedb}, and FoldX/ProteinSol/Aggrescan contracts)---closeable by formalizing each tool's algorithm in Lean~4
\item 20 ML-dependent predicates (ESMFold/AlphaFold structure~\cite{lin2023esm}, FoldX stability~\cite{schymkowitz2005foldx}, CamSol solubility~\cite{sormanni2015camsol}, NetMHCpan/MHCflurry immunogenicity~\cite{reynisson2020netmhcpan}, BePiPred, IEDB population contracts)---require novel research in ML model formalization
\end{itemize}

The BLOSUM62 conservation case (formerly the broad \texttt{blosum62\_tool\_sound} axiom and 2 \texttt{sorry}) has been fully discharged: the BLOSUM62 substitution matrix is now formalized in \texttt{BLOSUM62.lean} (11 theorems proving symmetry, diagonal non-negativity, and range bounds), and the \texttt{ConservationScore} case in \texttt{verification\_conditions\_imply\_property} is closed by \texttt{BLOSUM62.wellFormed\_proof} with no axiom. 0 \texttt{sorry} remain in \texttt{SLOTVerification.lean}.

\paragraph{Immunogenicity analysis is a separate API.} Immunogenicity analysis is available as a separate API (\texttt{compute\_immunogenicity()}, \texttt{deimmunize()}) that integrates NetMHCpan/MHCflurry when installed. It is \textbf{not} part of the default \texttt{optimize\_sequence()} path because it requires external tools and adds $\sim$seconds of latency. The \texttt{LowImmunogenicity} predicate is a SLOT predicate that returns \UNCERTAIN{} when tools are unavailable.

In CONSERVATIVE mode, SLOT predicates never return \PASS{} (vacuously sound by construction). In VERIFIED mode, they may return \PASS{} with evidence from the external tool, but the soundness proof for this mode has open obligations.

% ======================================================================
\section{Biosecurity Screening}\label{sec:biosecurity}
% ======================================================================

BioCompiler screens input sequences against 89 hazard signatures spanning 5 categories:

\begin{itemize}[leftmargin=*,itemsep=1pt]
\item Select agent toxins (ricin, botulinum, anthrax, shiga, diphtheria, tetanus, cholera, abrin)
\item Oncogenes (NRAS, KRAS, EGFR, VEGFA)
\item Viral surface proteins (SARS-CoV-2, HIV gp120, influenza HA)
\item Antibiotic resistance markers
\item Australia Group controlled pathogens (e.g., saxitoxin)
\end{itemize}

The screener uses exact matching (fast mode, $<$1\,ms) for pre-optimization screening, and fuzzy matching (Hamming + Levenshtein with k-mer pre-filtering) for post-optimization verification. Select-agent fuzzy matches at distance 2 are escalated to minimum risk ``medium'' to ensure detection.

All 8 reference hazard sequences are correctly flagged (8/8). The pre-optimization fast mode adds $<$1\,ms overhead, a feature DNAchisel lacks entirely (DNAchisel has no biosecurity screening).

% ======================================================================
\clearpage
\section{Evaluation}\label{sec:eval}
% ======================================================================

\subsection{Large-Scale Testing}

BioCompiler has been tested on 279 reference genes spanning 6 organisms (human, mouse, yeast, \textit{E.~coli}, CHO, \textit{Pichia}), totalling 1{,}674 protein--organism combinations across the bundled benchmark gene sets. Each combination runs the full pipeline: \texttt{optimize\_sequence()} $\to$ IR L0$\to$L1$\to$L2$\to$L3 $\to$ GenBank/FASTA/SBOL3 codegen, with protein preservation verification:

\begin{table}[!ht]
\centering
\begin{tabular}{lrrr}
\toprule
\textbf{Test Category} & \textbf{Combinations} & \textbf{Pass Rate} & \textbf{Failures} \\
\midrule
IR compiler (fast back-translation) & 100{,}068 & 100.0\% & 0 \\
E2E (integrated optimizer + IR + codegen) & 99{,}696 & 100.0\% & 0 \\
Head-to-head vs DNAchisel & 25 & 100.0\% & 0 \\
\midrule
\textbf{Total} & \textbf{199{,}789} & \textbf{100.0\%} & \textbf{0} \\
\bottomrule
\end{tabular}
\caption{Large-scale end-to-end test results.}
\label{tab:tests}
\end{table}

The reference gene set spans therapeutic proteins (insulin, growth hormone, erythropoietin, interferon $\alpha$2, IL-2, tPA, factor VIII, $\alpha$1-antitrypsin, glucocerebrosidase), \textit{E.~coli} benchmark genes (lacZ, lacI, recA, gyrA, rpoB, rpoD, dnaK, groEL, tufA, ampC, araC, malE, ompA), viral antigens (SARS-CoV-2 spike, SARS-CoV-2 RBD, influenza H1N1/H3N2 HA, HIV gp120, RSV F, rabies G, Zika E, dengue E, MTB Ag85B, \textit{P.~falciparum} CSP), and antibody targets (EGFR, HER2, VEGFA, PD-1). Stress-test sets cover extreme GC, extreme AT, 1-kb long sequences, and 10$\times$20 tandem repeats.

\subsection{Selenocysteine}

SECIS-aware selenocysteine translation is verified end-to-end on all 25 known human selenoproteins in \texttt{tests/test\_selenoproteins\_e2e.py} (130 tests). The catalog in \texttt{src/biocompiler/selenoproteins.py} records the verified Sec positions from the canonical selenoprotein literature (Kryukov et al.~2003; Gladyshev et al.~2016; UniProtKB feature tables). For each selenoprotein, the test constructs a fragment containing U at a verified position, back-translates to DNA with TGA at the Sec codon, runs the full L0$\to$L3 pipeline with SECIS annotation, and verifies that U is preserved. A negative control verifies that without SECIS annotation, the same TGA correctly translates to stop. The catalog includes the multi-Sec proteins SELENOP (10 Sec residues) and SELENOU (24 Sec residues).

\subsection{Retrospective Validation}

Retrospective validation on 24 documented therapeutic gene design failures (cryptic splicing, internal stops, CpG silencing, restriction sites, rare codon clusters) yields 79.2\% sensitivity (19/24 defects detected; 95\% CI: 57.9--92.1\% Wilson interval), with the threshold for cryptic promoters and the minimum restriction site length both adjusted to catch additional failure modes.

\subsection{Implementation Statistics}

\begin{table}[!ht]
\centering
\begin{tabular}{lr}
\toprule
\textbf{Component} & \textbf{Lines of Code} \\
\midrule
Python source & 149{,}230 \\
Lean~4 proofs & 8{,}746 \\
\quad Lean~4 theorems & 267 \\
Test suite & 162{,}169 \\
Standalone verifier & 473 \\
\midrule
Total & $\sim$321{,}000 \\
\bottomrule
\end{tabular}
\caption{Implementation statistics by component. Lean~4 theorem count is shown on a separate row because it is the headline verification metric (it is not additive with LOC).}
\label{tab:impl}
\end{table}

% ======================================================================
\section{Related Work}\label{sec:related}
% ======================================================================

\textbf{Proto BioCompiler} (Beal et al., PLOS ONE 2011)~\cite{beal2011protobiocompiler}---The name ``BioCompiler'' was first used by BBN Technologies for a Proto-to-genetic-regulatory-network compiler. BioCompiler (this work) is functionally different (codon optimization + type system + IR compiler) and we acknowledge the name collision.

\textbf{Cello} (Nielsen et al., Science 2016)~\cite{nielsen2016cello}---Verilog $\to$ DNA logic synthesis for transcriptional circuits. A different problem domain (circuit design, not codon optimization).

\textbf{DNAchisel} (Piret et al., 2020)~\cite{piret2020dnachisel}---A constraint-based DNA optimization library. On a 25-gene head-to-head benchmark, BioCompiler wins \emph{E.~coli} targets (higher CAI, faster) while DNAchisel wins human targets (higher CAI, faster); both achieve zero constraint violations. DNAchisel lacks formal verification, biosecurity screening, and an IR compiler.

\textbf{LinearDesign} (Zhang et al., Nature 2023)~\cite{zhang2023lineardesign}---mRNA vaccine design tool optimizing codon usage and mRNA secondary structure. A different problem domain (mRNA-level design, not codon-level verification).

\textbf{LLM-based circuit design tools} (2024--2026)---A growing body of work applies large language models to genetic circuit design, emitting SBOL or similar interchange formats. These tools represent the LLM-to-SBOL frontier and are complementary to BioCompiler's verified-codon-optimization approach: where LLM tools target high-level circuit specification, BioCompiler targets formally verified codon-level design with machine-checked lowering passes.

% ======================================================================
\section{Conclusion and Future Work}\label{sec:conclusion}
% ======================================================================

BioCompiler demonstrates that compiler-engineering principles---typed IRs, verified lowering passes, integrated constraint solving---can be successfully applied to gene design. The key results are:

\begin{enumerate}[leftmargin=*,itemsep=1pt]
\item A real five-level IR compiler with formally verified lowering correctness (Lean~4, 267 theorems, \texttt{compile\_correctness}).
\item An integrated optimizer that achieves zero constraint violations; on a 25-gene benchmark it wins \emph{E.~coli} (higher CAI, faster) and loses human (lower CAI, slower) to DNAchisel.
\item 1{,}674 test combinations across 279 reference genes and 6 organisms with a 100\% pass rate, plus 130 end-to-end tests covering all 25 known human selenoproteins (\texttt{tests/test\_selenoproteins\_e2e.py}) and 10 SECIS-aware unit tests (\texttt{tests/test\_secis.py}).
\item Biosecurity screening (8/8 hazards) and SECIS-aware selenocysteine support.
\end{enumerate}

\textbf{Limitations.} The formal verification covers 17 core deterministic predicates (the \emph{deterministic layer}---every predicate whose verdict can be decided by sequence computation alone). 0 \texttt{sorry} remain in \texttt{SLOTVerification.lean}; the 2 former BLOSUM62-related \texttt{sorry} have been discharged by formalizing the BLOSUM62 substitution matrix in \texttt{BLOSUM62.lean}. The 15 former broad tool-soundness \texttt{axiom} declarations have been narrowed to 34 specific, independently-testable contracts. \emph{Note: the axiom count increased from 15 to 34---the TCB is larger, but each assumption is narrower and independently testable at runtime via 34 evidence checks in \texttt{src/biocompiler/provenance/runtime\_evidence.py}.} These are \emph{tool-interface correctness properties} that cannot be discharged without formalizing the external tools themselves (ESMFold's neural network, FoldX's energy function, NetMHCpan's binding model), which is out of scope. The proof \emph{degrades gracefully}: in conservative mode the system never produces a false \PASS{} regardless. The system has not been validated in wet-lab or clinical settings. By analogy with CompCert~\cite{leroy2009compcert} and seL4~\cite{klein2014sel4}, the verified core is intended as a foundation on which larger verified gene-design stacks can be built; CakeML~\cite{kumar2014cakeml} is a similar model for end-to-end compiler verification.

\textbf{Future work.} (1) Discharge the 34 narrowed tool-soundness \texttt{axiom} declarations in \texttt{SLOTVerification.lean} by formalizing each external tool's algorithm in Lean~4 (the BLOSUM62 case is already discharged via \texttt{BLOSUM62.lean}; MaxEntScan~\cite{yeo2004maxentscan} is the next candidate---its splice-site scoring already uses the real Yeo \& Burge 2004 trained parameters, not approximations). (2) Wet-lab validation of optimized sequences. (3) Extend the IR to include UTR/polyA/m6A design (mRNA-level optimization).

% ======================================================================
\clearpage
\bibliographystyle{plain}
\bibliography{references}
% ======================================================================

\end{document}
